\documentclass[12pt, a4paper]{exam}
% Corrigé rédigé pour le site Ressources NSI et SNT (CC BY-SA 4.0).
% Compilation : pdflatex (deux passes).
\usepackage[T1]{fontenc}
\usepackage[utf8]{inputenc}
\usepackage{lmodern}
\usepackage[french]{babel}
\usepackage[a4paper,top=2.3cm,bottom=2.8cm,left=2cm,right=2cm,headsep=0.6cm,footskip=1.2cm]{geometry}
\usepackage{xcolor}
\usepackage{listings}
\usepackage{graphicx}
\usepackage{array}
\usepackage{amsmath}
\usepackage{amssymb}
\usepackage{enumitem}
\usepackage{needspace}
\usepackage{tikz}
\usetikzlibrary{positioning,arrows.meta,calc}
\usepackage[most]{tcolorbox}

% Coloration du code : même style que les sujets
\definecolor{codegreen}{rgb}{0,0.6,0}
\definecolor{codegray}{rgb}{0.5,0.5,0.5}
\definecolor{codepurple}{rgb}{0.58,0,0.82}
\definecolor{backcolour}{rgb}{0.95,0.95,0.92}
\lstdefinestyle{mystyle}{
    backgroundcolor=\color{backcolour},
    commentstyle=\color{codegreen}\itshape,
    keywordstyle=\color{magenta}\bfseries,
    numberstyle=\tiny\color{codegray},
    stringstyle=\color{codepurple},
    basicstyle=\footnotesize\ttfamily,
    breakatwhitespace=false,
    breaklines=true,
    postbreak=\mbox{\textcolor{codegray}{$\hookrightarrow$}\space},
    captionpos=b,
    keepspaces=true,
    columns=fullflexible,
    numbers=left,
    numbersep=5pt,
    xleftmargin=1.5em,
    showspaces=false,
    showstringspaces=false,
    showtabs=false,
    tabsize=4,
    upquote=true,
    extendedchars=true,
    literate={à}{{\`a}}1 {â}{{\^a}}1 {é}{{\'e}}1 {è}{{\`e}}1 {ê}{{\^e}}1 {ë}{{\"e}}1
             {î}{{\^i}}1 {ï}{{\"i}}1 {ô}{{\^o}}1 {ù}{{\`u}}1 {û}{{\^u}}1 {ç}{{\c c}}1
             {É}{{\'E}}1 {À}{{\`A}}1 {œ}{{\oe}}1 {…}{{\ldots}}1
}
\lstset{style=mystyle, language=Python}
\newcommand{\code}[1]{\texttt{#1}}

% Encadrés du corrigé (lisibles en noir et blanc)
\tcbset{corrbase/.style={enhanced, breakable, sharp corners, left=2mm, right=2mm, top=1mm, bottom=1mm,
  fonttitle=\bfseries\small, coltitle=black, before skip=2mm, after skip=3mm}}
\newtcolorbox{reponse}[1][Réponse]{corrbase, unbreakable, colback=gray!6, colframe=black, colbacktitle=gray!22, boxrule=0.8pt, title={#1}}
\newtcolorbox{erreur}[1][Erreur fréquente]{corrbase, colback=white, colframe=black, boxrule=0.4pt, leftrule=3pt, title={#1}, colbacktitle=white, attach title to upper={\ :\ }}
\newtcolorbox{retenir}[1][À retenir]{corrbase, colback=white, colframe=black, boxrule=0.4pt, borderline={0.4pt}{1.2pt}{black}, title={#1}, colbacktitle=white, attach title to upper={\ :\ }}

% En-tête d'une question : numéro, points, rappel de la question
\newcommand{\Q}[3]{\par\needspace{5\baselineskip}\bigskip\noindent{\large\bfseries Question #1}\ifx&#2&\else\ \textnormal{(#2)}\fi\par\nopagebreak\smallskip\noindent\textit{#3}\par\nopagebreak\smallskip}
\newcommand{\Qp}[2]{\par\medskip\noindent\textbf{#1}\ \textit{#2}\par\nopagebreak\smallskip}
\newcommand{\methode}{\par\noindent\textbf{Méthode.}\ }
\newcommand{\application}{\par\noindent\textbf{Application.}\ }

% Arbres
\tikzset{arbre/.style={every node/.style={draw,circle,minimum size=7mm,inner sep=1pt,font=\small},level distance=11mm, thick},
  marque/.style={fill=gray!35, line width=1.6pt},
  titre/.style={draw=none, rectangle, inner sep=2pt}}

\raggedbottom
\setlength{\parindent}{0pt}
\setlength{\parskip}{3pt}
\setlist{itemsep=1pt, topsep=2pt}
\newtcolorbox{remarque}[1][Remarque sur l'énoncé]{corrbase, colback=white, frame hidden, borderline={0.6pt}{0pt}{black, dashed}, title={#1}, colbacktitle=white, attach title to upper={\ :\ }}

\title{Corrigé --- Examen Term NSI - Tri, recherche et complexité}
\author{Babash - NSI Tani Malandi et Johanne Bakalara}
\date{20 points}

\begin{document}
\pagestyle{foot}
\footer{}{\footnotesize Babash - NSI Tani Malandi et Johanne Bakalara \textperiodcentered{} CC BY-SA 4.0 \textperiodcentered{} page \thepage}{}
\maketitle
\thispagestyle{foot}

\begin{retenir}[Comment lire ce corrigé]
Pour chaque question : la question rappelée en italique, la réponse attendue dans un cadre gris, puis l'explication (la méthode en étapes, puis son application). Les cadres « Erreur fréquente » montrent les pièges classiques. Essaie d'abord de refaire la question seul·e, puis compare.
\end{retenir}

\begin{retenir}[Le vocabulaire de la complexité]
\begin{itemize}
\item $n$ est la \textbf{taille des données} : le nombre d'éléments du tableau.
\item La \textbf{complexité} d'un algorithme dit comment son nombre d'opérations grandit quand $n$ grandit. On l'écrit avec un « grand O » : $O(n)$ se lit « grand O de $n$ ».
\item $O(n)$ (\textbf{linéaire}) : si $n$ est multiplié par 10, le temps est multiplié par 10 environ.
\item $O(n^2)$ (\textbf{quadratique}) : si $n$ est multiplié par 10, le temps est multiplié par 100.
\item $O(\log n)$ (\textbf{logarithmique}) : $\log_2 n$ est le nombre de fois qu'on peut diviser $n$ par 2 avant d'arriver à 1. Par exemple, $\log_2 1024 = 10$, car $2^{10} = 1024$.
\item Le \textbf{pire des cas} est la situation où l'algorithme fait le plus d'opérations ; le \textbf{meilleur des cas}, celle où il en fait le moins.
\end{itemize}
\end{retenir}

\Q{1 (a)}{2 points}{Classer par ordre croissant d'efficacité : $O(n^2)$, $O(1)$, $O(n \log n)$, $O(n)$, $O(2^n)$, $O(\log n)$.}
\begin{reponse}
Du moins efficace au plus efficace :
\[ O(2^n) \quad < \quad O(n^2) \quad < \quad O(n \log n) \quad < \quad O(n) \quad < \quad O(\log n) \quad < \quad O(1) \]
\end{reponse}
\methode
\begin{enumerate}
\item Un algorithme est d'autant plus \textbf{efficace} que son nombre d'opérations grandit \textbf{lentement} quand $n$ grandit.
\item Pour comparer les classes, calcule chacune pour une même valeur de $n$, par exemple $n = 1\,000$ (avec $\log_2 1\,000 \approx 10$).
\item Range-les de la plus grande valeur (la moins efficace) à la plus petite (la plus efficace).
\end{enumerate}
\application pour $n = 1\,000$ :
\begin{center}
\renewcommand{\arraystretch}{1.2}
\begin{tabular}{|l|l|r|}
\hline
\textbf{Classe} & \textbf{Nom} & \textbf{Opérations pour $n = 1\,000$} \\ \hline
$O(1)$ & constante & 1 \\
$O(\log n)$ & logarithmique & environ 10 \\
$O(n)$ & linéaire & 1\,000 \\
$O(n \log n)$ & quasi linéaire & environ 10\,000 \\
$O(n^2)$ & quadratique & 1\,000\,000 \\
$O(2^n)$ & exponentielle & un nombre de 302 chiffres \\ \hline
\end{tabular}
\end{center}
\begin{remarque}
« Ordre croissant d'efficacité » veut dire : du moins efficace au plus efficace. C'est l'ordre inverse de celui des complexités, qui va de $O(1)$ à $O(2^n)$. Une copie qui donne l'ordre $O(1) < O(\log n) < O(n) < O(n \log n) < O(n^2) < O(2^n)$ en précisant « de la plus efficace à la moins efficace » montre la même connaissance. Ce qui compte : dire clairement dans quel sens on range.
\end{remarque}
\begin{erreur}
Placer $O(n \log n)$ avant $O(n)$ en le croyant plus rapide. Pour $n \geq 2$, $n \log_2 n$ est plus grand que $n$ : $O(n \log n)$ est un peu moins efficace que $O(n)$.
\end{erreur}

\Q{1 (b)}{1 point}{Un tri fusion ($O(n \log n)$) met 1 ms pour 1\,000 éléments. Pour $10^6$ éléments, le temps devient-il « infini » ou reste-t-il de l'ordre de quelques secondes ?}
\begin{reponse}
Le temps reste de l'ordre de \textbf{quelques secondes} : environ \textbf{2 secondes}. \\
Les données sont 1\,000 fois plus nombreuses et $\log n$ est multiplié par 2 : le temps est multiplié par environ $1\,000 \times 2 = 2\,000$, donc $1 \text{ ms} \times 2\,000 = 2$ s.
\end{reponse}
\methode
\begin{enumerate}
\item On suppose que le temps est proportionnel à $n \log_2 n$.
\item On calcule par combien est multiplié $n \log_2 n$ quand on passe de $10^3$ à $10^6$.
\item On multiplie le temps de départ par ce nombre.
\end{enumerate}
\application
\begin{itemize}
\item $n$ passe de $10^3$ à $10^6$ : il est multiplié par $1\,000$.
\item $\log_2 n$ passe d'environ 10 ($2^{10} = 1\,024$) à environ 20 ($2^{20} \approx 1\,000\,000$) : il est multiplié par 2.
\item Donc $n \log_2 n$ est multiplié par $1\,000 \times 2 = 2\,000$, et le temps passe de 1 ms à $2\,000$ ms $= 2$ s.
\end{itemize}
\textbf{Pour comparer.} Un tri en $O(n^2)$ verrait son temps multiplié par $1\,000^2 = 1\,000\,000$ : il passerait d'1 ms à $1\,000$ s, soit plus d'un quart d'heure.

\Q{2}{2,5 points}{Relier chaque algorithme à la description de son principe.}
\begin{reponse}
\renewcommand{\arraystretch}{1.2}
\begin{tabular}{@{}l@{\ $\longrightarrow$\ }l@{}}
Recherche linéaire & Parcourir un par un tous les éléments. \\
Recherche dichotomique & Diviser l'espace de recherche par deux à chaque étape. \\
Tri par insertion & Insérer la valeur actuelle à sa place dans la partie triée. \\
Tri par sélection & Chercher le plus petit élément pour le placer au début. \\
Tri fusion & Diviser pour régner : trier deux moitiés puis les réunir. \\
\end{tabular}
\end{reponse}
\textbf{Les mots-clés à reconnaître.} « Linéaire » : on avance en ligne, un élément après l'autre. « Dichotomie » : couper en deux (du grec \emph{dikha}, en deux). « Insertion » : insérer. « Sélection » : choisir le minimum. « Fusion » : réunir deux listes triées en une seule.
\begin{erreur}
Confondre le tri par insertion et le tri par sélection. Repère : la \textbf{sélection} va \emph{chercher} le minimum dans la partie non triée ; l'\textbf{insertion} prend l'élément suivant tel qu'il vient et le \emph{glisse} à sa place dans la partie déjà triée.
\end{erreur}

\Q{3}{2,5 points}{Donner la classe de complexité de chaque algorithme.}
\begin{reponse}
Recherche linéaire : $O(n)$. \quad Recherche dichotomique : $O(\log n)$. \\
Tri par insertion : $O(n^2)$. \quad Tri par sélection : $O(n^2)$. \quad Tri fusion : $O(n \log n)$.
\end{reponse}
On donne la complexité dans le \textbf{pire des cas}, comme dans le cours :
\begin{itemize}
\item \textbf{Recherche linéaire} : au pire, on regarde les $n$ éléments un par un : $O(n)$.
\item \textbf{Recherche dichotomique} : à chaque étape, la zone de recherche est divisée par 2. On s'arrête au bout d'environ $\log_2 n$ étapes : $O(\log n)$.
\item \textbf{Tris par insertion et par sélection} : deux boucles imbriquées, qui font chacune jusqu'à $n$ tours, donc environ $n \times n$ opérations : $O(n^2)$.
\item \textbf{Tri fusion} : on coupe en deux environ $\log_2 n$ fois, et chaque « étage » de fusions traite les $n$ éléments : $O(n \log n)$.
\end{itemize}

\Q{4}{2,5 points}{Attribuer à chaque bloc de code (A à E) l'algorithme correspondant.}
\begin{reponse}
A : recherche dichotomique. \quad B : tri par insertion. \quad C : recherche linéaire. \\
D : tri par sélection. \quad E : tri fusion.
\end{reponse}
\methode Pour chaque code, cherche l'indice qui trahit l'algorithme :
\begin{itemize}
\item \textbf{Code A} : les bornes \code{g} et \code{d}, le milieu \code{m = (g + d) // 2}, puis \code{g = m + 1} ou \code{d = m - 1}. On garde une seule moitié : c'est la \textbf{dichotomie}.
\item \textbf{Code B} : on met de côté \code{k = t[i]}, puis la boucle \code{while} décale vers la droite (\code{t[j+1] = t[j]}) les éléments plus grands que \code{k}, et on pose \code{k} dans le trou. C'est l'\textbf{insertion}.
\item \textbf{Code C} : une seule boucle qui compare chaque \code{t[i]} à \code{v}. C'est la \textbf{recherche linéaire}.
\item \textbf{Code D} : la boucle intérieure cherche l'indice \code{m} du plus petit élément de \code{t[i:]}, puis on l'échange avec \code{t[i]}. C'est la \textbf{sélection}.
\item \textbf{Code E} : la fonction \code{tri} s'appelle elle-même sur les deux moitiés \code{t[:m]} et \code{t[m:]}, puis appelle \code{fusion}. C'est le \textbf{tri fusion}, un algorithme récursif.
\end{itemize}

\Q{5 (a)}{1 point}{Quelle condition le tableau \code{t} doit-il respecter pour que la recherche dichotomique fonctionne ?}
\begin{reponse}
Le tableau doit être \textbf{trié} (dans l'ordre croissant, pour le code A).
\end{reponse}
La dichotomie compare \code{v} au milieu, puis jette une moitié. Si \code{t[m] < v}, elle suppose que tout ce qui est à gauche de \code{m} est encore plus petit, donc inutile. Cette déduction n'est vraie que si le tableau est trié. \textbf{Contre-exemple.} Sur le tableau non trié \code{[12, 2, 18, 5, 8]}, cherchons 5. Le milieu est 18 : $18 > 5$, on garde la gauche (\code{[12, 2]}). Le milieu est 12 : $12 > 5$, on garde la gauche, qui est vide. L'algorithme conclut que 5 est absent, alors qu'il est à l'indice 3.

\Q{5 (b)}{1 point}{Recherche dichotomique avec \code{t = [2, 5, 8, 12, 18]} et \code{v = 12} : valeurs successives de \code{m}.}
\begin{reponse}
\code{m} vaut \textbf{2}, puis \textbf{3}. La valeur 12 est trouvée à l'indice 3.
\end{reponse}
\textbf{Tableau de trace} (code A) :
\begin{center}
\renewcommand{\arraystretch}{1.2}
\begin{tabular}{|c|c|c|c|c|l|}
\hline
\textbf{Tour} & \code{g} & \code{d} & \code{m = (g + d) // 2} & \code{t[m]} & \textbf{Décision} \\ \hline
1 & 0 & 4 & $(0 + 4) // 2 = 2$ & 8 & $8 < 12$ : \code{g = m + 1 = 3} \\
2 & 3 & 4 & $(3 + 4) // 2 = 3$ & 12 & $12 = 12$ : on renvoie 3 \\ \hline
\end{tabular}
\end{center}
\begin{erreur}
Donner les \emph{valeurs} \code{t[m]} (8 puis 12) au lieu des \emph{indices} \code{m} (2 puis 3). Autre piège : \code{//} est la division entière : $7 // 2 = 3$, pas 3,5.
\end{erreur}

\Q{5 (c)}{2 points}{Pour $n = 1\,000\,000$ éléments, nombre maximum de comparaisons de la recherche linéaire et de la recherche dichotomique (pire cas). Aide : $\log_2(250\,000) \approx 18$.}
\begin{reponse}
Recherche linéaire : \textbf{1\,000\,000} comparaisons. \\
Recherche dichotomique : environ \textbf{20} comparaisons (on compte un passage dans la boucle par comparaison avec \code{t[m]}).
\end{reponse}
\begin{itemize}
\item \textbf{Linéaire.} Au pire, la valeur est au dernier rang ou absente : on compare \code{v} aux $1\,000\,000$ éléments.
\item \textbf{Dichotomique.} À chaque tour, la zone est divisée par 2. On compte les divisions par 2 jusqu'à 1 élément : c'est $\log_2(1\,000\,000)$.
\end{itemize}
\textbf{Utiliser l'aide.} $1\,000\,000 = 250\,000 \times 4 = 250\,000 \times 2^2$. Donc $\log_2(1\,000\,000) = \log_2(250\,000) + 2 \approx 18 + 2 = 20$. Autrement dit : deux divisions par 2 font passer d'un million à $250\,000$, puis il en faut encore environ 18.

Une simulation du code A sur un tableau trié d'un million d'éléments donne au plus 20 tours de boucle. Un million de comparaisons contre 20 : c'est tout l'intérêt de trier les données avant de chercher.

\Q{6 (a)}{1,5 point}{Code B (tri par insertion) avec \code{t = [7, 2, 5, 3]} : état du tableau après chaque itération de la boucle \code{for i}.}
\begin{reponse}
Fin de l'itération $i = 1$ : \code{[2, 7, 5, 3]} \\
Fin de l'itération $i = 2$ : \code{[2, 5, 7, 3]} \\
Fin de l'itération $i = 3$ : \code{[2, 3, 5, 7]}
\end{reponse}
\textbf{Le code B, ligne par ligne.}
\begin{lstlisting}
for i in range(1, len(t)):
    k = t[i]
    j = i - 1
    while j >= 0 and t[j] > k:
        t[j+1] = t[j]
        j -= 1
    t[j+1] = k
\end{lstlisting}
\begin{itemize}
\item Ligne 1 : \code{i} va de 1 à \code{len(t) - 1}, ici 1, 2, 3. La partie \code{t[0..i-1]} est déjà triée.
\item Ligne 2 : \code{k} garde la valeur à insérer.
\item Lignes 4 à 6 : tant que l'élément de gauche est plus grand que \code{k}, on le décale d'une case vers la droite.
\item Ligne 7 : on pose \code{k} dans la case libérée.
\end{itemize}
\textbf{Tableau de trace} :
\begin{center}
\renewcommand{\arraystretch}{1.2}
\begin{tabular}{|c|c|l|c|}
\hline
\code{i} & \code{k} & \textbf{Ce que fait la boucle \code{while}} & \textbf{Tableau en fin de tour} \\ \hline
1 & 2 & $7 > 2$ : 7 décalé ; $j = -1$ : arrêt ; 2 posé en 0 & \code{[2, 7, 5, 3]} \\
2 & 5 & $7 > 5$ : 7 décalé ; $2 > 5$ faux : arrêt ; 5 posé en 1 & \code{[2, 5, 7, 3]} \\
3 & 3 & $7 > 3$, $5 > 3$ : décalés ; $2 > 3$ faux ; 3 posé en 1 & \code{[2, 3, 5, 7]} \\ \hline
\end{tabular}
\end{center}
\begin{erreur}
Commencer à $i = 0$, ou ajouter une ligne $i = 4$ : \code{range(1, len(t))} s'arrête à \code{len(t) - 1}, soit 3. Autre erreur : confondre avec la sélection, qui donnerait \code{[2, 3, 5, 7]} dès la fin du deuxième tour. L'insertion ne va pas chercher le minimum : elle traite les éléments dans l'ordre où ils arrivent.
\end{erreur}

\Q{6 (b)}{1 point}{Quelle proposition est l'invariant de la boucle \code{for} ?}
\begin{reponse}
La deuxième : « À l'étape $i$, la tranche \code{t[0...i]} contient les mêmes éléments qu'au départ, mais triés. »
\end{reponse}
Un \textbf{invariant de boucle} est une propriété vraie à chaque tour de boucle. Elle sert à prouver que l'algorithme est correct : à la fin, $i = n - 1$, et tout le tableau est trié. On teste chaque proposition sur la trace de la question 6 (a) :
\begin{itemize}
\item Proposition 1 : fausse. Après $i = 1$, le tableau est \code{[2, 7, 5, 3]} : le dernier élément, 3, n'est pas le plus grand. Cette propriété décrit d'autres tris, comme le tri à bulles.
\item Proposition 2 : vraie. Après $i = 1$, \code{[2, 7]} est trié ; après $i = 2$, \code{[2, 5, 7]} ; après $i = 3$, \code{[2, 3, 5, 7]}. Ce sont bien les éléments de départ, réordonnés.
\item Proposition 3 : fausse. Après $i = 2$, \code{t[2]} vaut 7 ; mais à la fin, 7 est à l'indice 3 : il n'était pas à sa place définitive. Cette propriété correspond au tri par sélection.
\end{itemize}

\Q{6 (c) i}{1 point}{Meilleur des cas pour le tri par insertion, et sa complexité.}
\begin{reponse}
Le meilleur des cas est un tableau \textbf{déjà trié} (dans l'ordre croissant). La complexité est alors \textbf{linéaire}, en $O(n)$.
\end{reponse}
Si le tableau est trié, la condition \code{t[j] > k} est fausse dès le premier test : la boucle \code{while} ne fait aucun tour. On fait seulement une comparaison par tour de \code{for}, soit $n - 1$ comparaisons. Pour $n = 1\,000$, une simulation donne bien 999 comparaisons.

\Q{6 (c) ii}{1 point}{Pire des cas pour le tri par insertion, et sa complexité.}
\begin{reponse}
Le pire des cas est un tableau trié dans l'\textbf{ordre décroissant}. La complexité est alors \textbf{quadratique}, en $O(n^2)$.
\end{reponse}
Chaque nouvel élément est plus petit que tous ceux de la partie triée : il faut le décaler jusqu'au début. Le tour $i$ fait $i$ décalages, d'où $1 + 2 + \cdots + (n-1) = \frac{n(n-1)}{2}$ comparaisons. Pour $n = 1\,000$ : $499\,500$, ce que confirme la simulation. C'est de l'ordre de $n^2$.

\Q{6 (c) iii}{1 point}{Pourquoi dit-on que le tri par insertion est « efficace » sur une liste presque triée, par rapport au tri par sélection ?}
\begin{reponse}
Sur une liste presque triée, le tri par insertion fait très peu de décalages : la boucle \code{while} s'arrête presque tout de suite, et le coût est proche de $n$. Le tri par sélection, lui, parcourt toujours toute la partie non triée pour chercher le minimum : il fait toujours $\frac{n(n-1)}{2}$ comparaisons, même si la liste est déjà triée.
\end{reponse}
\textbf{Exemple chiffré.} Pour 1\,000 éléments triés, sauf deux paires d'éléments voisins échangées, une simulation donne 1\,001 comparaisons pour l'insertion, contre 499\,500 pour la sélection : environ 500 fois moins.
\begin{erreur}
Répondre seulement « parce qu'il est plus rapide ». La justification attendue explique \emph{pourquoi} : la boucle intérieure de l'insertion s'arrête dès qu'elle trouve la place, alors que celle de la sélection va toujours jusqu'au bout du tableau.
\end{erreur}

\end{document}
